\chapter{Đam Mê Trải Nghiệm}
\markboth{Đam Mê Trải Nghiệm}{Đam Mê Trải Nghiệm}

\section{Cái Gì Cũng Muốn Thử Một Chút}
\begin{truyen}{Cái Gì Cũng Muốn Thử Một Chút}{Wanting to Try a Little of Everything}
\chuhoa{H}{ọc sinh:} ``Em không muốn bị đóng khung. Em muốn trải nghiệm nhiều để tìm ra bản thân.''
\textit{(Student: ``I do not want to be boxed in. I want many experiences so I can discover myself.'')}

\teacherpair{Trải nghiệm là tốt nếu nó làm em sâu hơn. Nhưng nhiều bạn bây giờ nghiện cảm giác mới lạ hơn là nghiện học tới nơi tới chốn. Cái gì cũng thử một ít nghe có vẻ sống động, nhưng nếu không bám được vào thứ gì đủ lâu, em sẽ chỉ thu về một bộ sưu tập bề mặt.}{``Experience is good if it deepens you. But many students today are addicted more to novelty than to following anything through. Trying a little of everything sounds lively, but if you cannot stay with anything long enough, you end up collecting surfaces, not substance.''}
\end{truyen}

\section{Đi Nhiều Không Đồng Nghĩa Là Lớn Nhanh}
\begin{truyen}{Đi Nhiều Không Đồng Nghĩa Là Lớn Nhanh}{Going Many Places Does Not Mean Growing Faster}
\chuhoa{H}{ọc sinh:} ``Em thấy bạn bè em đi workshop, trại hè, câu lạc bộ đủ kiểu. Em sợ mình đứng yên.''
\textit{(Student: ``My friends attend workshops, camps, and clubs of every kind. I am afraid of standing still.'')}

\teacherpair{Đứng yên không phải lúc nào cũng thua. Có khi người khác đang chạy vòng ngoài còn em đang đào sâu. Trải nghiệm nhiều không xấu. Nhưng nếu em để nỗi sợ bị tụt trend kéo mình đi, em sẽ rất khó biết đâu là thứ đáng ở lại đủ lâu để biến thành năng lực thật.}{``Standing still is not always losing. Sometimes others are running around the surface while you are digging deep. Many experiences are not bad. But if fear of missing trends drags you around, you will struggle to know what deserves your long enough attention to become real skill.''}
\end{truyen}

\section{Trải Nghiệm Tốt Phải Để Lại Một Vết Tích}
\begin{truyen}{Trải Nghiệm Tốt Phải Để Lại Một Vết Tích}{A Good Experience Leaves a Mark}
\chuhoa{H}{ọc sinh:} ``Vậy trải nghiệm sao cho đáng?''
\textit{(Student: ``Then how do I make experiences worthwhile?'')}

\teacherpair{Sau mỗi trải nghiệm, em phải thấy mình đổi một chút: biết hơn, chịu hơn, sâu hơn, hoặc rõ mình hơn. Nếu đi rất nhiều mà tính khí, kỷ luật và chiều sâu vẫn y nguyên, thì phần lớn em chỉ đang tiêu thụ trải nghiệm như người ta tiêu thụ nội dung thôi.}{``After each experience, you should change a little: know more, endure more, deepen more, or understand yourself more clearly. If you go through many things while your character, discipline, and depth remain unchanged, then you are merely consuming experiences the way people consume content.''}
\end{truyen}

\begin{baihoc}
Trải nghiệm chỉ thật sự quý khi nó được tiêu hóa thành phẩm chất.
\textit{(Experience becomes valuable only when digested into character.)}
\end{baihoc}

\begin{ghinhoanh}
\textbf{Short English to remember:}

Experience is not decoration.

It should change your depth.
\end{ghinhoanh}

\begin{tuhocnhanh}
1. Em đang tìm trải nghiệm để lớn lên, hay để bớt sợ bị bỏ lại? \textit{(Are you seeking experience to grow, or to fear being left behind less?)}

2. Trải nghiệm gần nhất nào đã thực sự để lại một thay đổi trong em? \textit{(Which recent experience truly left a change in you?)}
\end{tuhocnhanh}

\ngancach
